\documentclass[11pt]{article}
% The preamble every project's paper shares: % The preamble every project's paper shares: % The preamble every project's paper shares: \input{../../preamble} from
% projects/<project>/paper/. Boards come from chessboard and xskak, plots from
% pgfplots; every number a paper prints comes from its generated.tex.
\usepackage[LSBC4,T1]{fontenc}
\usepackage[utf8]{inputenc}
\usepackage{lmodern}
\usepackage[margin=1in]{geometry}
\usepackage{microtype}
\usepackage{amsmath}
\usepackage{booktabs}
\usepackage{tabularx}
\usepackage{array}
\usepackage{xcolor}
\usepackage[obeyspaces,spaces]{url}
\usepackage{xskak}
\usepackage{pgfplots}
\pgfplotsset{compat=1.18}
\usepackage{caption}
\usepackage{subcaption}
\usepackage{enumitem}
\usepackage[section]{placeins}
\usepackage[hidelinks]{hyperref}

\setchessboard{boardfontencoding=LSBC4, boardfontfamily=skaknew, blackfieldcolor=black!22,
  whitefieldcolor=white, setfontcolors, showmover=false, label=true, labelfontsize=6pt,
  boardfontsize=14pt}

% A move or line in UCI, and a position in FEN, both breakable at a slash.
\newcommand{\uci}[1]{\texttt{#1}}
\newcommand{\fen}[1]{{\footnotesize\path{#1}}}
\newcolumntype{L}{>{\raggedright\arraybackslash}X}

\setlength{\parskip}{0.4em}
\setlength{\parindent}{0pt}
 from
% projects/<project>/paper/. Boards come from chessboard and xskak, plots from
% pgfplots; every number a paper prints comes from its generated.tex.
\usepackage[LSBC4,T1]{fontenc}
\usepackage[utf8]{inputenc}
\usepackage{lmodern}
\usepackage[margin=1in]{geometry}
\usepackage{microtype}
\usepackage{amsmath}
\usepackage{booktabs}
\usepackage{tabularx}
\usepackage{array}
\usepackage{xcolor}
\usepackage[obeyspaces,spaces]{url}
\usepackage{xskak}
\usepackage{pgfplots}
\pgfplotsset{compat=1.18}
\usepackage{caption}
\usepackage{subcaption}
\usepackage{enumitem}
\usepackage[section]{placeins}
\usepackage[hidelinks]{hyperref}

\setchessboard{boardfontencoding=LSBC4, boardfontfamily=skaknew, blackfieldcolor=black!22,
  whitefieldcolor=white, setfontcolors, showmover=false, label=true, labelfontsize=6pt,
  boardfontsize=14pt}

% A move or line in UCI, and a position in FEN, both breakable at a slash.
\newcommand{\uci}[1]{\texttt{#1}}
\newcommand{\fen}[1]{{\footnotesize\path{#1}}}
\newcolumntype{L}{>{\raggedright\arraybackslash}X}

\setlength{\parskip}{0.4em}
\setlength{\parindent}{0pt}
 from
% projects/<project>/paper/. Boards come from chessboard and xskak, plots from
% pgfplots; every number a paper prints comes from its generated.tex.
\usepackage[LSBC4,T1]{fontenc}
\usepackage[utf8]{inputenc}
\usepackage{lmodern}
\usepackage[margin=1in]{geometry}
\usepackage{microtype}
\usepackage{amsmath}
\usepackage{booktabs}
\usepackage{tabularx}
\usepackage{array}
\usepackage{xcolor}
\usepackage[obeyspaces,spaces]{url}
\usepackage{xskak}
\usepackage{pgfplots}
\pgfplotsset{compat=1.18}
\usepackage{caption}
\usepackage{subcaption}
\usepackage{enumitem}
\usepackage[section]{placeins}
\usepackage[hidelinks]{hyperref}

\setchessboard{boardfontencoding=LSBC4, boardfontfamily=skaknew, blackfieldcolor=black!22,
  whitefieldcolor=white, setfontcolors, showmover=false, label=true, labelfontsize=6pt,
  boardfontsize=14pt}

% A move or line in UCI, and a position in FEN, both breakable at a slash.
\newcommand{\uci}[1]{\texttt{#1}}
\newcommand{\fen}[1]{{\footnotesize\path{#1}}}
\newcolumntype{L}{>{\raggedright\arraybackslash}X}

\setlength{\parskip}{0.4em}
\setlength{\parindent}{0pt}

% Written by projects/focus-backtest/paper/generate.mjs; do not edit by hand.

\newcommand{\AppCommit}{d8b771238c}
\newcommand{\Window}{40}
\newcommand{\Baseline}{100}
\newcommand{\SwitchRatio}{1.5}
\newcommand{\SuccessP}{0.9}
\newcommand{\ArchiveGames}{634}
\newcommand{\Weeks}{ten}
\newcommand{\Focuses}{eight}
\newcommand{\Finished}{seven}
\newcommand{\Met}{five}
\newcommand{\NotMet}{two}
\newcommand{\Switches}{eight}
\newcommand{\Graded}{nine}
\newcommand{\Filled}{eight}
\newcommand{\GamesPerWeek}{50}
\newcommand{\ChanceOfMet}{14}
\newcommand{\MeanReal}{40}
\newcommand{\Pairs}{90}
\newcommand{\StillThere}{40}
\newcommand{\BrierReal}{0.246}
\newcommand{\BrierBase}{0.240}
\newcommand{\WeakestReal}{0.50}
\newcommand{\WeakestGap}{0.03}
\newcommand{\WeakestSe}{3.19}
\newcommand{\Runs}{4,000}
\newcommand{\Pool}{634}
\newcommand{\Perms}{400}
\newcommand{\Level}{10}
\newcommand{\TableForty}{%
\begin{tabular}{lrrrr}
\toprule
grade, 40 games & no change & 10\% cut & 25\% cut & 50\% cut \\
\midrule
the app & 44.3 & 55.3 & 71.1 & 87.6 \\
\quad its CUSUM alarm alone & 35.9 & 45.4 & 59.3 & 74.4 \\
\quad its posterior alone & 28.1 & 39.5 & 60.8 & 84.8 \\
fresh (per game) & 8.3 & 15.5 & 31.1 & 63.2 \\
fresh, per move & 9.1 & 20.0 & 44.1 & 76.3 \\
random cause, the app's grade & 35.5 & & & \\
\bottomrule
\end{tabular}
}
\newcommand{\TableEighty}{%
\begin{tabular}{lrrrr}
\toprule
grade, 80 games & no change & 10\% cut & 25\% cut & 50\% cut \\
\midrule
the app & 51.9 & 65.0 & 80.6 & 94.0 \\
\quad its CUSUM alarm alone & 44.8 & 55.9 & 69.7 & 83.7 \\
\quad its posterior alone & 31.4 & 48.4 & 72.8 & 92.5 \\
fresh (per game) & 10.0 & 19.0 & 41.3 & 74.0 \\
fresh, per move & 9.8 & 23.5 & 54.9 & 82.4 \\
random cause, the app's grade & 41.8 & & & \\
\bottomrule
\end{tabular}
}
\newcommand{\TableHundred}{%
\begin{tabular}{lrrrr}
\toprule
grade, 120 games & no change & 10\% cut & 25\% cut & 50\% cut \\
\midrule
the app & 55.0 & 69.1 & 84.0 & 95.2 \\
\quad its CUSUM alarm alone & 46.0 & 57.5 & 72.0 & 86.3 \\
\quad its posterior alone & 35.4 & 54.6 & 77.9 & 94.8 \\
fresh (per game) & 10.0 & 20.4 & 46.7 & 78.2 \\
fresh, per move & 9.6 & 25.6 & 60.7 & 84.8 \\
random cause, the app's grade & 43.6 & & & \\
\bottomrule
\end{tabular}
}
\newcommand{\NullApp}{44}
\newcommand{\NullAppLong}{55}
\newcommand{\NullCusum}{36}
\newcommand{\NullPosterior}{28}
\newcommand{\NullControl}{36}
\newcommand{\NullControlPosterior}{19}
\newcommand{\NullFresh}{8.3}
\newcommand{\NullPerMove}{9.1}
\newcommand{\NotMetUnclear}{87}
\newcommand{\QuarterApp}{71}
\newcommand{\QuarterPerMove}{44}
\newcommand{\QuarterPerMoveLong}{61}
\newcommand{\HalfPerMove}{76}
\newcommand{\HalfPerMoveLong}{85}
\newcommand{\QuarterAppNet}{27}
\newcommand{\QuarterPerMoveNet}{35}
\newcommand{\TargetCut}{8}
\newcommand{\TargetCutLo}{2}
\newcommand{\TargetCutHi}{38}
\newcommand{\NeededCut}{31}
\newcommand{\Cv}{1.5}
\newcommand{\DriftWindow}{10}
\newcommand{\DriftOlder}{9}
\newcommand{\TopCause}{endgame.tense}
\newcommand{\TopCauseShare}{31}


\title{Does the weekly focus's grade tell a real improvement from none?\\[0.3em]
\large A replay of the ONE focus on the player's games, and a simulation with a known truth}
\author{Chess Coach research lab}
\date{}

\begin{document}
\maketitle

\begin{abstract}
Every Sunday the weekly review picks one focus, the leak cause where the player gives away the
most against his opponents, freezes a baseline and a target, and grades it over the next \Window{}
games. A met focus is the plan's evidence that the training worked. Replayed on his archive, the
picker changed focus \Switches{} times in \Weeks{} weeks and graded \Met{} of \Finished{} finished
focuses met. In \Runs{} simulated focuses drawn from his own games with nothing changed, the grade
says met \NullApp\% of the time at \Window{} games and \NullAppLong\% at 120, against the \Level\% its
own \SuccessP{} threshold implies; a cause drawn at random gets \NullControl\%, so most of it is the test, not the choice. A
permutation test against older games that played no part in the choice holds its level
(\NullFresh\% and \NullPerMove\%), and reading points per move of the cause sees a real cut
sooner than points per game. The targets ask for a median cut of \TargetCut\% while \Window{}
games can only see about \NeededCut\%. A candidate grade built on those findings says met
\CandNullLong\% of the time with nothing changed and \CandQuarterLong\% with a quarter of the loss
removed, over \CandWindow{} games; on his games it graded not met a focus the app had graded met
on a baseline read at its worst.
\end{abstract}

\section{The question}

The weekly's ONE focus is the app's answer to ``what should he work on now''. It ranks the leak
causes (an opening, middlegame or endgame position that is tense or quiet, under a minute on the
clock, a fast move in a tense position, the move after a gift, late in a game) by the Elo they
cost, the probability the leak is real, its relevance, its actionability and its freshness, and
keeps a running focus unless a challenger scores \SwitchRatio{} times as much. A chosen focus is
pre-registered: a baseline (his points lost on the cause per game over the last \Baseline{}
games), a target halfway to his opponents' level, and a CUSUM test. It is graded over the next
\Window{} games: met when the CUSUM alarms downward or the posterior probability that his loss
fell reaches \SuccessP, else not met.

Two things decide whether that grade can be read as evidence: how often it says met when
nothing changed, and how large a change it needs before it says met. Both are measured here on
the app's own code at commit \texttt{\AppCommit}, through the lab's bridge.

\section{What the app did with his games}

The replay runs the weekly's own choice and commit every Sunday of the archive
(\ArchiveGames{} games), on a metrics instance that holds only the games played by then. He plays
a median of \GamesPerWeek{} games a week, more than the \Window-game window, so in \Filled{} of the
\Graded{} weeks that had a focus running, it had already been decided by the next Sunday, and the
picker moved on: \Focuses{} focuses in \Weeks{} weeks. Of the \Finished{} that finished, \Met{}
were graded met and \NotMet{} not met. At the simulated null rate below, \Met{} or more of
\Finished{} would read met by chance alone \ChanceOfMet\% of the time.

The picker's probability that a leak is real averaged \MeanReal\% over the \Pairs{} ranked leaks
of the replayed weeks, and \StillThere\% of them were still there (his loss on the cause still
above his opponents') over the next \Window{} games. It did not say which: its Brier score for
that outcome was \BrierReal, against \BrierBase{} for a constant. A 40-game future is noisy, so
this does not show the probability is wrong; it shows that, over the span a focus runs, it does
not tell the leaks that will show up again from those that will not. One switch went to a leak
the picker gave even odds (\WeakestReal), with a gap of \WeakestGap{} points a game against a
standard error of \WeakestSe.

\section{A simulation with a known truth}

Each of \Runs{} runs draws, with replacement from his \Pool{} archived games, an older
\Baseline-game history, the \Baseline{} games the picker reads, and a window of up to 120 games
after it, so nothing changes between them. The app's \texttt{rankFocus}, \texttt{preregister} and
\texttt{gradeFocus} run on the draws exactly as the weekly would. The same draws are then graded
with a known share of his loss on the chosen cause removed: each of his moves on the cause loses
nothing with that probability, from a fixed uniform per move, so every cut is compared on the
same games. Beside the app's grade, four candidates were scored on the same draws (fixed before
this run, after a first run measured the app's grade alone): the app's grade on a cause drawn at
random from the ranked ones (no selection), the app's grade without its CUSUM pass, and a
one-sided permutation test (\Perms{} permutations, level \Level\%) of the window against the
\emph{older} games, which played no part in choosing the cause, on points per game and on points
per move of the cause.

\begin{table}[h]
\centering
\TableForty

\medskip
\TableEighty

\medskip
\TableHundred
\caption{The share of simulated focuses graded met, by grade, window and the share of his loss
on the cause removed. With no change, every share above the level (\Level\%) is a false met.}
\end{table}

\paragraph{With nothing changed.} The app's grade says met \NullApp\% of the time at \Window{}
games, and more as the window grows (\NullAppLong\% at 120): the CUSUM alone fires
\NullCusum\% of the time and the posterior alone \NullPosterior\%. On a cause drawn at random the
grade still says met \NullControl\% of the time (the posterior alone \NullControlPosterior\%), so
selection adds little; the test itself is the problem. His per-game loss on a cause is heavily
skewed (a median coefficient of variation of \Cv): most games lose little and a few lose a lot,
so a short window usually sits below the mean, and a CUSUM tuned to a normal-looking shift
alarms on it. The permutation tests, which make no such assumption, hold their level.

\paragraph{The choice inflates the baseline.} With nothing changed, the window's loss on the
chosen cause sits a median \DriftWindow\% below the baseline that chose it, and the older games
sit \DriftOlder\% below it too: the picker reads the baseline at its worst. A test against the
baseline credits that regression to the mean to the training. The older games carry no such
bias.

\paragraph{When something changed.} With a quarter of his loss on the cause removed, the app's
grade says met \QuarterApp\% of the time at \Window{} games, but \NullApp\% of that it would say
anyway: net \QuarterAppNet{} points. Points per move against the older games say met
\QuarterPerMove\% of the time (net \QuarterPerMoveNet{} points), \QuarterPerMoveLong\% at 120
games; with half of it removed, \HalfPerMove\% and \HalfPerMoveLong\%. Per move beats per game at
every cut and window, because a game with many moves on the cause and a game with one are no
longer weighed alike.

\paragraph{Targets the window can see.} The pre-registered target asks for a median cut of
\TargetCut\% of his loss on the cause (from \TargetCutLo\% at the tenth percentile to
\TargetCutHi\% at the ninetieth). For the posterior to reach \SuccessP{} on \Window{} games, the cut has to be
about \NeededCut\% at the median spread. Most targets are below what the window can tell from
noise.

\paragraph{Not met is mostly ``cannot tell''.} Of the focuses graded not met with nothing
changed, \NotMetUnclear\% had a posterior between 0.1 and 0.9: the grade had no evidence either
way, but the weekly records it as not met.

\section{What it means for the app}

The grade is the evidence the plan reads, so it should earn its word. The simulation points to
five changes, each measurable again on this harness before it ships:

\begin{enumerate}[nosep]
\item Grade a focus against games that did not choose it (the games before the baseline, or a
held-out half), not against the baseline that picked it.
\item Drop the CUSUM early pass, or calibrate its threshold by resampling his own games, so a
skewed loss does not alarm by itself.
\item Read points per move of the cause, not per game.
\item Say three things, not two: improved, not improved, and not enough games yet.
\item Size the window, or hold the focus, to the cut the target asks for: at his pace a focus
can run two or three weeks of games instead of one.
\end{enumerate}

\section{A candidate grade}

The candidate (\texttt{candidates/permove/focus.ts}) is the app's \texttt{focus.ts} with the first
four changes and a longer hold. It grades a focus by a one-sided permutation test, \CandPerms{}
shuffles of games seeded by the focus, of his points per move on the cause in the window against
the \Baseline{} games before the ones the choice read (the baseline itself when the archive is
shorter). It says met when the test's p-value is at most \MetP, not met when it is
\NotMetP{} or more (the window is no better than the older games), and unclear in between, and it
holds a focus for \CandWindow{} games with no early pass. It ran through the same simulation on the
same draws, again on fresh draws (seed \FreshSeed), and through the same weekly replay.

\paragraph{It holds its level.} With nothing changed, it says met \CandNull\% of the time at 40
games, \CandNullMid\% at 80 and \CandNullLong\% at \CandWindow, and \CandNullFresh\% on fresh draws;
on a cause drawn at random, \CandControl\%. On the fresh draws the app's grade says met
\AppNullFresh\% of the time at 40 games, so its false met is not a property of one set of draws.

\paragraph{It sees a real cut.} With a tenth, a quarter and a half of his loss on the cause
removed, it says met \CandTenthLong\%, \CandQuarterLong\% and \CandHalfLong\% of the time at
\CandWindow{} games (\CandQuarter\% at 40 for a quarter). These are the per-move test's rates
above, within two points at every cut and window: the candidate is that test, with its own
shuffles.

\paragraph{What the weekly would say.} At the end of a \CandWindow-game window with nothing
changed, the candidate says met \VNullMet\%, unclear \VNullUnclear\% and not met \VNullNotMet\% of the
time; with a quarter of the loss removed, \VQuarterMet\%, \VQuarterUnclear\% and \VQuarterNotMet\%.
Not met now means something: it is common when nothing changed and rare when something did.

\paragraph{On his games.} Replayed over the same \Weeks{} weeks, the candidate held \CandFocuses{}
focuses instead of \Focuses, each for \CandSpanMin{} to \CandSpanMax{} weeks. It graded
\CandUnclear{} of them unclear and \CandNotMet{} not met, and none met. The one it graded not met,
\texttt{\FlipCause}, is one the app graded met: his loss on the cause fell from \FlipGameBaseline{}
points a game in the baseline to \FlipGameAppWindow{} in the \Window{} games after it, but the
\Baseline{} games before the baseline had read \FlipGameReference. Per move, the windows after the
choice (\FlipAppWindow{} over \Window{} games, \FlipWindow{} over \CandWindow) were no better than
the older games' \FlipReference, while the baseline that chose it read \FlipBaseline. The app's met
was the regression to the mean the simulation predicts.

\section{Limits}

The simulation draws games independently, so it has no streaks, no rating drift and no change in
his opponents; a real improvement may arrive unevenly. The cut removes loss uniformly across his
moves on the cause, which is one shape of improvement among many. The replay covers \Weeks{}
weeks, too few to grade the picker's choices on outcomes; it shows how the loop behaves, not
whether it helped. The candidate compares a window with games played weeks earlier, so anything
else that changed in between (his openings, his opponents, his rating) reads as a change too.

\end{document}
